\chapter[Global Comparison]{Global Comparison with Classical Statistical Divergences}
\section{Introduction}
Chapter 25 established that the pullback metric induced by the Universal
Optimality Defect (UOD) is locally equivalent to the Fisher information
metric. The present chapter investigates the global relationship between
the UOD and the principal statistical divergences.
Throughout,
\[ \mathcal{D}(P,Q)=\|F(P)-F(Q)\|^{2}, \qquad
a=\log(P/Q). \]
Only rigorously established comparisons are presented.
\section{Local versus Global Behaviour}
The UOD is globally defined through Euclidean distance in logarithmic
quotient coordinates. Consequently:
\begin{itemize}
  \item locally it agrees with Fisher geometry to second order;
  \item globally it measures logarithmic distortions rather than additive
\end{itemize}
    probability differences.
\section{Comparison with Total Variation}
The Total Variation distance is
$TV(P,Q)=\frac12\sum_{i}\|P_i-Q_i\|.$
\begin{theorem}
\label{ch:thm:26-1}
There exist sequences of probability distributions for which the Total
Variation distance remains bounded while the UOD diverges to infinity.
\textbf{Proof.}
Let
$Q=\left(\frac13,\frac13,\frac13\right),$
and
\[ P_\varepsilon= \left( \varepsilon,
\frac{1-\varepsilon}{2}, \frac{1-\varepsilon}{2}
\right). \]
Then
\[ TV(P_\varepsilon,Q)\longrightarrow\frac23,
\]
whereas
$\log\frac{\varepsilon}{1/3}\longrightarrow-\infty.$
After projection onto the quotient hyperplane, the logarithmic diameter
still diverges. By Theorem 23.5,
$\mathcal{D}(P_\varepsilon,Q)\longrightarrow\infty.$
Hence no bounded function of Total Variation can globally characterize
the UOD. 
\end{theorem}
\section{Comparison with Hellinger Distance}
The squared Hellinger distance is
\[ H^{2}(P,Q)= \frac12
\sum_{i}(\sqrt{P_i}-\sqrt{Q_i})^{2}. \]
\begin{theorem}
\label{ch:thm:26-2}
The Hellinger distance is bounded whereas the UOD is unbounded.
\textbf{Proof.}
Using
$H^{2}(P,Q) = 1-\sum_{i}\sqrt{P_iQ_i},$
one immediately obtains
$0\le H(P,Q)\le1.$
By Theorem 26.1, the UOD diverges along suitable sequences approaching
the boundary of the simplex. Therefore the UOD cannot be represented
globally as a function of the Hellinger distance. 
\end{theorem}
\section{Local Comparison with Pearson's Chi-Squared Divergence}
Pearson's divergence is
\[ \chi^{2}(P,Q) = \sum_{i}
\frac{(P_i-Q_i)^2}{Q_i}. \]
Let
\[ P=Q+\delta, \qquad
\sum_{i}\delta_{i}=0. \]
Using
$\log(1+x)=x-\frac{x^2}{2}+O(x^{3}),$
one obtains
\[ \chi^{2}(P,Q) = \sum_{i}
\frac{\delta_i^2}{Q_i} + O(\|\delta\|^{3}), \]
whereas
\[ \mathcal{D}(P,Q) = \sum_{i}
\frac{\delta_i^2}{Q_i^2} - \frac1n \left(
\sum_{i}\frac{\delta_i}{Q_i} \right)^{2} +
O(\|\delta\|^{3}). \]
Thus both quantities possess quadratic local expansions, but they induce
different quadratic forms.
\section{First Structural Comparison with KL}
The Kullback--Leibler divergence satisfies
\[ D_{KL}(P\|Q) = \sum_{i} P_i
\log\frac{P_i}{Q_i}. \]
Introducing the logarithmic field
$a_i=\log\frac{P_i}{Q_i},$
we obtain the identity
\[ \boxed{
D_{KL}(P\|Q)=E_P[a].
} \]
By Chapter 21,
\[ \boxed{
\mathcal{D}(P,Q)
=
n\,Var_U(a),
} \]
where (U) denotes the uniform distribution.
Hence the KL divergence and the UOD are both functionals of the same
log-likelihood field, but evaluated in fundamentally different ways: the
KL is its expectation with respect to (P), whereas the UOD is its
centered second moment with respect to the uniform measure.
This observation motivates the functional framework developed in the
next chapter.
\section{Summary}
The rigorous comparison established in this chapter shows:
\begin{itemize}
  \item the UOD is locally equivalent to Fisher geometry;
  \item it is globally distinct from Total Variation;
  \item it is globally distinct from Hellinger distance;
  \item it shares only the local quadratic regime with Pearson's divergence;
  \item together with KL, it naturally arises from the same logarithmic
\end{itemize}
    field, although through different functionals.
These results conclude the comparison with classical divergences and
prepare the functional theory developed in Part V.
